\documentclass[xcolor=dvipsnames, xcolor=table]{beamer} % Classe de documento para apresentações
\usepackage{amsmath, amssymb}
\usepackage{listings}
\usepackage{listings}\usetheme{Berlin}
\usecolortheme{beaver}
\usepackage{textpos}
\usepackage{color}
\usepackage{wasysym}
\renewcommand\tabcolsep{4pt}
\usepackage{csquotes} 
\usepackage{wrapfig}
\usepackage{microtype}
\usepackage[labelfont=bf]{caption}
\usepackage{epigraph}
\usepackage[english]{babel}
\usepackage{float}
\usepackage{hyperref}
\usepackage{amsmath}
\newcommand{\dd}[1]{\mathrm{d}#1}
\usepackage{amsfonts}
\usepackage{amssymb}
\usepackage{amsbsy}
\usepackage{graphicx}
\usepackage{mathtools}

\usepackage{booktabs}

\usepackage{adjustbox}
\hypersetup{pdfpagemode=FullScreen} % para o pdf abrir automaticamente em modo full screen
\setbeamertemplate{caption}[numbered] % para numerar as tabelas
\usepackage{subcaption}
\usepackage[none]{hyphenat} 
\usepackage{bm} %%% bold vectors
\usepackage{tabularx}  
\usepackage{tikz}




\setbeamercolor{author in head/foot}{fg=Maroon}     % cor do nome do autor no rodape
\setbeamercolor{institute in head/foot}{fg=Maroon}  % cor do nome do ISEG no rodape
\setbeamercolor{frametitle}{fg=Maroon, bg=black!9} % cor da letra e do fundo fa faixa com o nome de cada slide
\setbeamercolor{caption name}{fg=Maroon, bg=black!9}
\setbeamercolor*{title}{fg=white, bg=Maroon}          % cor do titulo e da caixa de titulo na capa
\setbeamercolor{itemize item}{fg=Maroon}               % cor dos bullets do ambiente \itemize 


\makeatletter % criaçao de um ambiente sem TOC no cabeçalho do slide
\newenvironment{noheadline}{
\setbeamertemplate{headline}{}
\addtobeamertemplate{frametitle}{\vspace*{-0.9\baselineskip}}{}
}{}
\makeatother

\usepackage[beamer,customcolors]{hf-tikz}

\tikzset{hl/.style={
    set fill color=red!80!black!40,
    set border color=red!80!black,
  },
}


\title[Matrices]{Matrices}
\author[João Vieira \& Pedro Fonseca]{\textbf {João Vieira \& Pedro Fonseca}}
\titlegraphic{\includegraphics[width=1.5cm]{yes}}
\institute[Introduction to R Programming]{\textbf {Introduction to R Programming}}
\date{\today}


% position the yes
\addtobeamertemplate{frametitle}{}{%
\begin{textblock*}{100mm}(.85\textwidth,-.9cm)
\includegraphics[scale=.05]{yes}
\end{textblock*}}



%%%%%%%%%%%%% Ambiente de teorema/definição com as cores do ISEG %%%%%%

\makeatletter
\def\th@mystyle{%
    \normalfont % body font
    \setbeamercolor{block title example}{bg=Maroon,fg=white}
    \setbeamercolor{block body example}{bg=black!9,fg=black}
    \def\inserttheoremblockenv{exampleblock}
  }
\makeatother
\theoremstyle{mystyle}
\newtheorem*{defi}{Definição}

\setbeamertemplate{defis}[numbered]



%%%%%%%%%%%%%%%% Bibliografia %%%%%%%%%%%%%%%%%%%%%%%

\usepackage{url}
\usepackage[style=authoryear, bibencoding=utf8, minnames=1, maxnames=3,
maxbibnames=99,backref=true, natbib=true, dashed=false, terseinits=true, 
firstinits=true, uniquename=false, uniquelist=true, labeldate=true, 
doi=false, isbn=false, natbib=true, backend=biber]{biblatex}
\DefineBibliographyStrings{english}{%
    backrefpage = {Cited on page},
    backrefpages = {Cited on pages},
}
% Change the default formatting to be more "statistical"
\DeclareFieldFormat{url}{\url{#1}}
\DeclareFieldFormat[article]{pages}{#1}
\DeclareFieldFormat[inproceedings]{pages}{\lowercase{pp.}#1}
\DeclareFieldFormat[incollection]{pages}{\lowercase{pp.}#1}
\DeclareFieldFormat[article]{volume}{\mkbibbold{#1}}
\DeclareFieldFormat[article]{number}{\mkbibparens{#1}}
\DeclareFieldFormat[article]{title}{\MakeCapital{#1}}
\DeclareFieldFormat[article]{url}{}
\DeclareFieldFormat[book]{url}{}
\DeclareFieldFormat[inbook]{url}{}
\DeclareFieldFormat[incollection]{url}{}
\DeclareFieldFormat[inproceedings]{url}{}
\DeclareFieldFormat[inproceedings]{title}{#1}
\DeclareFieldFormat{shorthandwidth}{#1}
% No dot before number of articles
\usepackage{xpatch}
\xpatchbibmacro{volume+number+eid}{\setunit*{\adddot}}{}{}{}
% Remove In: for an article.
\renewbibmacro{in:}{%
  \ifentrytype{article}{}{%
  \printtext{\bibstring{in}\intitlepunct}}}
% Get rid of months in citations
\AtEveryBibitem{\clearfield{month}}
\AtEveryCitekey{\clearfield{month}}
\setlength{\parindent}{1,3cm}
\raggedbottom

\renewcommand\bibfont{\scriptsize}
% If you have more than one page of references, you want to tell beamer
% to put the continuation section label from the second slide onwards
\setbeamertemplate{frametitle continuation}[from second]
% Now get rid of all the colours
\setbeamercolor*{bibliography entry title}{fg=black}
\setbeamercolor*{bibliography entry author}{fg=black}
\setbeamercolor*{bibliography entry location}{fg=black}
\setbeamercolor*{bibliography entry note}{fg=black}





\begin{document}
%%%%%%%%%%%%%%%%%%%%%  CAPA  %%%%%%%%%%%%%%%%%%%%%%%%%%%%%
\begin{noheadline}

\begin{frame}%[plain]
\vfill
\centering

\begin{beamercolorbox}[sep=8pt,center,colsep=-4bp,rounded=true,shadow=true]{title}
\usebeamerfont{title}\inserttitle\par%
\ifx\insertsubtitle\@empty%
\else%
\vskip0.25em%
{\usebeamerfont{subtitle}\usebeamercolor[fg]{subtitle}\insertsubtitle\par}%
\fi%     
\end{beamercolorbox}%

\vskip1em\par

\begin{beamercolorbox}[sep=8pt,center,colsep=-4bp,rounded=true,shadow=true]{author}
\usebeamerfont{author}\insertauthor
\end{beamercolorbox}

{\usebeamercolor[fg]{titlegraphic}\inserttitlegraphic\par}

\begin{beamercolorbox}[sep=8pt,center,colsep=-4bp,rounded=true,shadow=true]{institute}
\usebeamerfont{institute}\insertinstitute
\end{beamercolorbox}

\begin{beamercolorbox}[sep=8pt,center,colsep=-4bp,rounded=true,shadow=true]{date}
\usebeamerfont{date}\insertdate
\end{beamercolorbox}\vskip0.5em

\end{frame}
\end{noheadline}



%%%%%%%%%%%%%% ITEMIZE SHAPES AND COLLERS%%%%%%%%%%%%%%%%%%%%%%%%%

\setbeamertemplate{itemize item}{\color{Maroon}\newmoon}
\setbeamertemplate{itemize subitem}{\color{Maroon}$\blacktriangleright$}



%%%%%%%%%%%%%%%%%%%%%%%%%%%%%%%%%%%%%%%%%%%%%%%%%%%%%%%%%%%%%%%

\section{Definition}

\begin{frame}[fragile] 
\frametitle{What is a Matrix?}

\begin{itemize} 
\item A matrix with $m$ rows and $n$ columns is defined as:
\end{itemize}

\begin{equation*}
A_{m,n} = 
\begin{pmatrix}
a_{1,1} & a_{1,2} & \cdots & a_{1,n} \\
a_{2,1} & a_{2,2} & \cdots & a_{2,n} \\
\vdots  & \vdots  & \ddots & \vdots  \\
a_{m,1} & a_{m,2} & \cdots & a_{m,n} 
\end{pmatrix}
\end{equation*}

\end{frame}

\section{Creating a Matrix}

\begin{frame}[fragile]
\frametitle{The \textit{matrix} function}

\begin{itemize}
\item To create a matrix use the \textit{matrix} function:

\begin{verbatim}
> A <- matrix(data, nrow, ncol, byrow)
\end{verbatim}

\item Only the \textit{data} argument is mandatory.
\item Provide the entries of the matrix as a vector to the data argument. 
\item With \textit{nrow} you can provide the number of row and with \textit{ncol} the number of columns.
\item By default, \textit{matrix} returns a $\operatorname{lenght}(\operatorname{data}) \times 1$ matrix. 

\end{itemize}

\end{frame}

\begin{frame}[fragile] 
\frametitle{The \textit{matrix} function: examples}

\begin{verbatim}
> A <- matrix(data = c(1, 2, 3 ,4 ,5 , 6), nrow = 2, 
ncol = 3)
> A
     [,1] [,2] [,3]
[1,]    1    3    5
[2,]    2    4    6

> b <- c(1, 5 ,5, 7)  # Remark: R is case sensitive!
> B <- matrix(b, nrow = 2, ncol = 2)
> B
     [,1] [,2]
[1,]    1    5
[2,]    5    7
\end{verbatim}



\end{frame}


\begin{frame}[fragile]
\frametitle{The \textit{byrow} argument}

\begin{itemize}
\item \textit{byrow} is a logical argument (TRUE or FALSE), set to FALSE by default.
\item If FALSE the output matrix is filled by columns, otherwise the matrix is filled by rows.
\end{itemize}

\end{frame}

\begin{frame}[fragile]
\frametitle{The \textit{byrow} argument: example}

\begin{verbatim}
> A <- matrix(data = c(1, 2, 3 ,4 ,5 , 6), nrow = 2, 
ncol = 3)
> A
     [,1] [,2] [,3]
[1,]    1    3    5
[2,]    2    4    6

> B <- matrix(data = c(1, 2, 3 ,4 ,5 , 6), nrow = 2, 
ncol = 3, byrow = TRUE)
> B
     [,1] [,2] [,3]
[1,]    1    2    3
[2,]    4    5    6
\end{verbatim}
\end{frame}



\begin{frame}[fragile] %%%%%%%%%%%%%%%%%%%%%%% FRAME %%%%%%%%%%%%%%%%%%%%%%%%%%%%
\frametitle{Argument Names}

\begin{itemize}
\item Remember: Argument names are optional:
\begin{verbatim}
> matrix(data = c(1, 2, 3 ,4 ,5 , 6), nrow = 2, 
ncol = 3)

     [,1] [,2] [,3]
[1,]    1    3    5
[2,]    2    4    6
 
> matrix(c(1, 2, 3 ,4 ,5 , 6),  2,  3)

     [,1] [,2] [,3]
[1,]    1    3    5
[2,]    2    4    6
 
\end{verbatim}
\end{itemize}

\end{frame}


\begin{frame}[fragile] %%%%%%%%%%%%%%%%%%%%%%% FRAME 
\frametitle{Recycling}

\begin{itemize}
\item R \textit{recycles} the \textit{data} vector if its length is not compatible with the number of elements of the output matrix:

\begin{verbatim}
> matrix(1:5, nrow = 2, ncol = 3, byrow = TRUE)
     [,1] [,2] [,3]
[1,]    1    2    3
[2,]    4    5    1
\end{verbatim}
\end{itemize}


\end{frame}



\begin{frame}[fragile] %%%%%%%%%%%%%%%%%%%%%%% FRAME %%%%%%%%%%%%%%%%%%%%%%%%%%%%
\frametitle{\textit{rbind} and \textit{cbind}}

\begin{itemize}
\item If you have multiple vectors of equal length, you create matrix by binding them together with the \textit{rbind} and \textit{cbind} functions:

\begin{verbatim}
> rbind(1:3, 4:6)
     [,1] [,2] [,3]
[1,]    1    2    3
[2,]    4    5    6

\end{verbatim}
\end{itemize}

\end{frame}

\begin{frame}[fragile]
\frametitle{\textit{rbind} and \textit{cbind}: example}

\begin{verbatim}

> set.seed(123)
 
> a <- 1:4
> b <- seq(from = 1, by = 0.5, length.out = 4)
> c <- rpois(4, lambda = 2)
 
> rbind(a, b, c)
  [,1] [,2] [,3] [,4]
a    1  2.0    3  4.0
b    1  1.5    2  2.5
c    1  3.0    2  4.0

\end{verbatim}

\end{frame}

\begin{frame}[fragile]
\frametitle{\textit{rbind} and \textit{cbind}: example}

\begin{verbatim}
> mymatrix <- cbind(c(1,4), c(2,5), c(3,6))
> mymatrix

     [,1] [,2] [,3]
[1,]    1    2    3
[2,]    4    5    6
\end{verbatim}

\end{frame}
\begin{frame}[fragile] %%%%%%%%%%%%%%%%%%%%%%% FRAME %%%%%%%%%%%%%%%%%%%%%%%%%%%%

\frametitle{Dimensions}

\begin{itemize}
\item Number of columns:

\begin{verbatim}
> ncol(mymatrix)
[1] 3
\end{verbatim}


\item Number of rows:

\begin{verbatim}
> nrow(mymatrix)
[1] 2
\end{verbatim}

\item Both:

\begin{verbatim}
> dim(mymatrix)
[1] 2 3

> dim(mymatrix)[1]
[1] 2

\end{verbatim}

\end{itemize}

\end{frame}

\begin{frame}[fragile]

\begin{itemize}
\item Matrix rows and columns can be named:

\begin{verbatim}
> M <- matrix(c(1, 5, 3, 7), ncol = 2, byrow = TRUE)

> rownames(M) <- c("x1", "x2")
> colnames(M) <- c("y1", "y2")

> M
   y1 y2
x1   1  5
x2   3  7
\end{verbatim}

\end{itemize}

\end{frame}

\section{Subsetting}

\begin{frame}[fragile] %%%%%%%%%%%%%%%%%%%%%%% FRAME %%%%%%%%%%%%%%%%%%%%%%%%%%%%

\frametitle{Subsetting}

\begin{itemize}
\item The first column:

\begin{verbatim}
> mymatrix[, 1]
[1] 1 4
\end{verbatim}


\item The first row:

\begin{verbatim}
> mymatrix[1, ]
[1] 1 2 3
\end{verbatim}

\item A specific element:

\begin{verbatim}
> mymatrix[2, 3]
[1] 6
\end{verbatim}

\end{itemize}


\end{frame}



\begin{frame}[fragile] %%%%%%%%%%%%%%%%%%%%%%% FRAME %%%%%%%%%%%%%%%%%%%%%%%%%%%%
\frametitle{Subsetting}

\begin{itemize}
\item The first and third columns:

\begin{verbatim}
> mymatrix[, c(1,3)]
     [,1] [,2]
[1,]    1    3
[2,]    4    6
\end{verbatim}


\item The first and second rows:

\begin{verbatim}
> mymatrix[c(1, 2), ]
     [,1] [,2] [,3]
[1,]    1    2    3
[2,]    4    5    6
\end{verbatim}

\item The elements in the diagonal of the matrix:

\begin{verbatim}
> diag(mymatrix)
[1] 1 5

\end{verbatim}

\end{itemize}

\end{frame}



\begin{frame}[fragile] %%%%%%%%%%%%%%%%%%%%%%% FRAME %%%%%%%%%%%%%%%%%%%%%%%%%%%%
\frametitle{Subsetting}

\begin{itemize}
\item To omit elements from a matrix, use negative indexes:
\begin{itemize}

\item The first column:

\begin{verbatim}
> mymatrix[, -1]
     [,1] [,2]
[1,]    2    3
[2,]    5    6
\end{verbatim}

\item The first row:

\begin{verbatim}
> mymatrix[-1, ]
[1] 4 5 6
\end{verbatim}

\item A specific row and column:

\begin{verbatim}
> mymatrix[-2, -3]
[1] 1 2
\end{verbatim}

\end{itemize}
\end{itemize}

\end{frame}


\begin{frame}[fragile] %%%%%%%%%%%%%%%%%%%%%%% FRAME %%%%%%%%%%%%%%%%%%%%%%%%%%%%
\frametitle{Subsetting}

\begin{itemize}
\item To omit elements from a matrix, use negative indexes:
\begin{itemize}
\item The first and third columns:

\begin{verbatim}
> mymatrix[, -c(1,3)]
[1] 2 5
\end{verbatim}


\item The first and third rows:

\begin{verbatim}
> mymatrix[-c(1,2), ]
     [,1] [,2] [,3]
\end{verbatim}

\end{itemize}
\end{itemize}



\end{frame}

\begin{frame}[fragile]
\frametitle{Subsetting}

\begin{itemize}
\item In a matrix with named columns or vectors, we can use the names to subset:

\begin{verbatim}
> M
   y1 y2
x1  1  5
x2  3  7

> M[, "y1"]
x1 x2 
 1  3 
 
> M["x2", "y1"]
[1] 3
\end{verbatim}

\end{itemize}

\end{frame}


\begin{frame}[fragile] %%%%%%%%%%%%%%%%%%%%%%% FRAME %%%%%%%%%%%%%%%%%%%%%%%%%%%%

\frametitle{Substituting}

\begin{itemize}
\item Overwriting elements of a matrix:

\begin{itemize}


\item The first column:

\begin{verbatim}
> mymatrix

     [,1] [,2] [,3]
[1,]    1    2    3
[2,]    4    5    6

> mymatrix[, 1] <- c(5, 5)
> mymatrix

     [,1] [,2] [,3]
[1,]    5    2    3
[2,]    5    5    6

\end{verbatim}

\end{itemize}

\end{itemize}

\end{frame}

\begin{frame}[fragile]
\frametitle{Substituting}

\begin{itemize}
\item Overwriting elements of a matrix:

\begin{itemize}

\item The first row:

\begin{verbatim}

> mymatrix

     [,1] [,2] [,3]
[1,]    1    2    3
[2,]    4    5    6

> mymatrix[1,] <- 30:32
> mymatrix
     [,1] [,2] [,3]
[1,]   30   31   32
[2,]    5    5    6
\end{verbatim}

\end{itemize}

\end{itemize}

\end{frame}


\begin{frame}[fragile] %%%%%%%%%%%%%%%%%%%%%%% FRAME %%%%%%%%%%%%%%%%%%%%%%%%%%%%
\frametitle{Substituting}

\begin{itemize}
\item Overwriting a specific element in a matrix:

\begin{verbatim}
> mymatrix[2,3] <- 60
> mymatrix
     [,1] [,2] [,3]
[1,]   30   31   32
[2,]    5    5   60
\end{verbatim}

\item Overwriting the diagonal:
\begin{verbatim}
> diag(mymatrix) <- rep(x = 0, times = 2)
> mymatrix
     [,1] [,2] [,3]
[1,]    0   31   32
[2,]    5    0   60
\end{verbatim}

\end{itemize}

\end{frame}

\section{Operations with Matrices}

\begin{frame}[fragile] %%%%%%%%%%%%%%%%%%%%%%% FRAME %%%%%%%%%%%%%%%%%%%%%%%%%%%%
\frametitle{Summation, subtraction and multiplication of scalars}

\begin{verbatim}
> A
     [,1] [,2]
[1,]    1    3
[2,]    2    4

> 2*A
     [,1] [,2]
[1,]    2    6
[2,]    4    8

> A + 5
     [,1] [,2]
[1,]    6    8
[2,]    7    9
\end{verbatim}

\end{frame}

\begin{frame}[fragile] %%%%%%%%%%%%%%%%%%%%%%% FRAME %%%%%%%%%%%%%%%%%%%%%%%%%%%%

\frametitle{Adding and subtracting matrices}

\begin{verbatim}
> A <- matrix(1:4, nrow = 2)
> B <- matrix(rep(1, times = 4), ncol = 2)

> A + B
     [,1] [,2]
[1,]    2    4
[2,]    3    5

> A - B
     [,1] [,2]
[1,]    0    2
[2,]    1    3

\end{verbatim}

\end{frame}

\begin{frame}[fragile] %%%%%%%%%%%%%%%%%%%%%%% FRAME %%%%%%%%%%%%%%%%%%%%%%%%%%%%

\frametitle{Element-wise multiplication and division}

\begin{verbatim}
> A <- matrix(1:4, nrow = 2)
> C <- matrix(rep(2, times = 4), ncol = 2)

> A * C     # This is not matrix multiplication!
     [,1] [,2]
[1,]    2    6
[2,]    4    8

> A / C
     [,1] [,2]
[1,]  0.5  1.5
[2,]  1.0  2.0

\end{verbatim}

\end{frame}

\section{Matrix Algebra}

\begin{frame}[fragile] %%%%%%%%%%%%%%%%%%%%%%% FRAME %%%%%%%%%%%%%%%%%%%%%%%%%%%%

\frametitle{Matrix Algebra}

\begin{itemize}
\item Transpose matrix:

\begin{verbatim}
> A
     [,1] [,2]
[1,]    1    3
[2,]    2    4

> t(A)
     [,1] [,2]
[1,]    1    2
[2,]    3    4
\end{verbatim}

\end{itemize}

\end{frame}

\begin{frame}[fragile] %%%%%%%%%%%%%%%%%%%%%%% FRAME %%%%%%%%%%%%%%%%%%%%%%%%%%%%

\frametitle{Matrix Algebra}

\begin{itemize}

\item Inverse matrix:

\begin{verbatim}
> A
     [,1] [,2]
[1,]    1    3
[2,]    2    4

> solve(A)
     [,1] [,2]
[1,]   -2  1.5
[2,]    1 -0.5
\end{verbatim}

\end{itemize}

\end{frame}

\begin{frame}[fragile] %%%%%%%%%%%%%%%%%%%%%%% FRAME %%%%%%%%%%%%%%%%%%%%%%%%%%%%

\frametitle{Matrix Algebra}

\begin{verbatim}

> A <- rbind(c(1, 3), c(2, 4))
> B <- matrix(1:4, ncol = 2, byrow = TRUE)

> A
     [,1] [,2]
[1,]    1    3
[2,]    2    4

> B
     [,1] [,2]
[1,]    1    2
[2,]    3    4
\end{verbatim}

\end{frame}

\begin{frame}[fragile] %%%%%%%%%%%%%%%%%%%%%%% FRAME %%%%%%%%%%%%%%%%%%%%%%%%%%%%

\frametitle{Matrix Algebra}

\begin{itemize}
\item Determinant of a matrix
\end{itemize}

\begin{verbatim}
> A <- matrix(c(-1, 1, 2, 
                -5, 3, 4,
                10, 8, -4), 
              ncol = 3, 
              byrow = TRUE)
              
> det(A)
[1] -76
\end{verbatim}

\end{frame}

\begin{frame}[fragile] %%%%%%%%%%%%%%%%%%%%%%% FRAME %%%%%%%%%%%%%%%%%%%%%%%%%%%%

\frametitle{Matrix Algebra}

\begin{itemize}

\item Matrix multiplication:

\begin{verbatim}
> A %*% B
     [,1] [,2]
[1,]   10   14
[2,]   14   20
\end{verbatim}

\item Matrix multiplication is not commutative:

\begin{verbatim}
> B %*% A
     [,1] [,2]
[1,]    5   11
[2,]   11   25
\end{verbatim}

\end{itemize}

\end{frame}

\begin{frame}[fragile]
\frametitle{Systems of Linear Equations}

\begin{itemize}
\item The \textit{solve} function can also be used to solve systems of linear equations.
\item \textit{solve(A, B)} computes $X$ from equation $A X = B$
\begin{itemize}
\item A: matrix with the coefficients of the equations
\item B: vector or matrix of the equation's right side
\item X: vector or matrix of unknowns
\end{itemize}
\item The default value of $B$ is an identity matrix. This is why \textit{solve}(A) gives $\operatorname{A}^{-1}$:
\end{itemize}

\begin{equation}
AX = \operatorname{I} \Leftrightarrow X = A^{-1}
\end{equation}

\end{frame}

\begin{frame}[fragile]
\frametitle{Systems of Linear Equations}

\begin{itemize}
\item Example 1:

\begin{itemize}
\item $5x = 10$, what is $x$?

\begin{verbatim}
> solve(5,10)
[1] 2
\end{verbatim}

\end{itemize}

\end{itemize}

\end{frame}

\begin{frame}[fragile]
\frametitle{Systems of Linear Equations}

\begin{itemize}

\item Example 2: $\begin{cases} 3x+2y & = 8  \\ \, \, \, x+  \, \, \,y & = 2  \end{cases}$ \, \,  what is $(x, y)$?

\begin{verbatim}
> A <- matrix(c(3,1,2,1), nrow = 2, ncol = 2)
> A
     [,1] [,2]
[1,]    3    2
[2,]    1    1

> b  <- c(8,  2)

> solve(A,b)
[1]  4 -2
\end{verbatim}

\item Conclusion: $x = 4$ and $y = -2$.

\end{itemize}

\end{frame}

\begin{frame}[fragile]
\frametitle{Systems of Linear Equations}

\begin{itemize}
\item Alternative solution to example 2:

\begin{itemize}
\item $Ax = b \Leftrightarrow x =  A^{-1}b $

\begin{verbatim}

> solve(A,b) 
[1]  4 -2

> solve(A) %*% b
     [,1]
[1,]    4
[2,]   -2

\end{verbatim}
\end{itemize}
\end{itemize}

\end{frame}

\begin{frame}[fragile] %%%%%%%%%%%%%%%%%%%%%%% FRAME %%%%%%%%%%%%%%%%%%%%%%%%%%%%
\frametitle{Questions?}


\begin{flushleft}“One whose knowledge is confined to books and whose wealth is in the possession of others, can use neither his knowledge nor wealth when the need for them arises.”\end{flushleft}\begin{flushright}
― Chanakya
\end{flushright}


\end{frame}




\end{document}


